% QMB_n10_spin_qubit.tex — Chapter 10: Quantum Dot, Spin, Stability
% Substantive basis: Doc. 173 (March 2026), Doc. 174, Doc. 178; own presentation
\chapter{Quantum Dot, Spin, and the Limit of Superposition}
\label{ch:spin}

\section{What a Quantum Dot Is: Geometry Enforces Discreteness}

A quantum dot is an artificial atom --- a cavity so small
that its geometry enforces discrete energy levels. Not the number
of atoms determines the levels, but the shape of the cavity: how the
tones of a flute depend on the length of the tube, not on the number
of atoms in the metal. The energy levels of the confined particle:
\begin{equation}
  E_n = \frac{n^2\pi^2\hbar^2}{2m_{\mathrm{eff}}R^2},\qquad n = 1,2,3,\dots
\end{equation}
The factor $\pi^2$ follows from the geometry: A standing wave that
vanishes at the boundary enforces $k\cdot R = \pi$ for the fundamental mode. The
discreteness comes from the geometry, not from a postulate.

In the FFGFT, an electron is not a point charge in a box, but
a stable torsion configuration of the underlying field. The
discrete levels are the allowed geometric configurations of this
field twist in a bounded space. Between the levels there is
no stable configuration; the field exists only in defined
patterns, not in arbitrary intermediate states.

The fundamental parameter $\xi = 4/30000$ generates a universal
length hierarchy $L_N = L_0\cdot(1/\xi)^N$ starting at the
fundamental length $L_0 = \xi\ell_P$. Level $N = 8$ corresponds to
$L_8\approx22$\,nm; the excitonic Bohr radius for CdSe lies
at $N\approx7{,}85$ --- the same scale level. The scale of excitons
is, in the FFGFT, not an empirical quantity, but geometrically
enforced.\status{S}

\begin{center}
\small
\begin{tabular}{cll}
\toprule
$N$ & $L_N$ & Physical Domain \\
\midrule
0 & $2{,}15\times10^{-39}$\,m & Fundamental FFGFT length \\
1 & $1{,}62\times10^{-35}$\,m & Planck length \\
6 & $3{,}8\times10^{-16}$\,m & Subnuclear \\
7 & $2{,}9\times10^{-12}$\,m & Atomic scale \\
\textbf{8} & $\mathbf{2{,}2\times10^{-8}}$\,\textbf{m} & Quantum dot and exciton range \\
9 & $1{,}6\times10^{-4}$\,m & Cellular scale \\
10 & $\approx1$\,m & Macroscopic \\
\bottomrule
\end{tabular}
\end{center}

\section{Two Conditions for Quantum Dot Behavior}

Two conditions must both be fulfilled:

\textbf{Condition 1 --- Geometric perfection} (absolute, temperature-
independent): The resonator must be precise enough to enforce standing waves. A defective lattice shows no confinement
quantum behavior, no matter how small, no matter at what temperature. Condition 1
is the harder one; it cannot be replaced by cooling.

\textbf{Condition 2 --- Level spacing versus heat} (size- and
temperature-dependent): The spacing $\Delta E = 3E_1$ between the lowest
two levels must exceed the thermal energy: $\Delta E > k_BT$.

From Condition 2 follows the maximum resonator size:
\begin{equation}
  R_{\mathrm{bulk}}(T) = \pi\hbar\sqrt{\frac{3}{2m_{\mathrm{eff}}k_BT}}.
\end{equation}
For silicon at 300\,K: $R_{\mathrm{bulk}}\approx15$\,nm --- exactly where
classical transistors stopped scaling in 2012.

\begin{center}
\small
\begin{tabular}{lrrrrr}
\toprule
\textbf{Material} & $m_{\mathrm{eff}}/m_e$ & \multicolumn{4}{c}{$R_{\mathrm{bulk}}$ at} \\
\cmidrule(lr){3-6}
 & & 4\,K & 77\,K & 300\,K & 310\,K \\
\midrule
GaAs & 0{,}067 & 221\,nm & 50\,nm & 25{,}5\,nm & 25{,}1\,nm \\
CdSe & 0{,}130 & 159\,nm & 36\,nm & 18{,}3\,nm & 18{,}0\,nm \\
Silicon & 0{,}190 & 131\,nm & 30\,nm & 15{,}2\,nm & 14{,}9\,nm \\
\bottomrule
\end{tabular}
\end{center}

Cooling to 4\,K shifts the limit for silicon from 15\,nm to
131\,nm --- a factor~9. This makes macroscopically acting objects into
quantum objects: A 50-nm silicon crystal is classical at room temperature,
quantized at 4\,K. A macroscopic body does not have
no levels, but $10^{23}$ so densely packed that $k_BT$ covers them all.
The discreteness does not disappear; it becomes so fine-grained
that it appears as band structure.

\section{Transistors and the $\xi$ Level Table}

When the semiconductor industry had to switch to FinFET in 2012, because
classical planar transistors had stopped scaling, the
channel lengths had fallen below $\xi$ level $N = 8{,}00$:

\begin{center}
\small
\begin{tabular}{p{3.5cm}ccp{2.8cm}}
\toprule
\textbf{Technology} & \textbf{Channel Length} & \textbf{$N$-Level} & \textbf{Regime} \\
\midrule
65\,nm (2005) & 65\,nm & 8{,}12 & classical \\
22\,nm FinFET (2012) & 22\,nm & \textbf{8{,}00} & quantum influences begin \\
10\,nm (2016) & 10\,nm & 7{,}91 & tunneling measurable \\
5\,nm (2020) & 5\,nm & 7{,}84 & QD regime \\
3\,nm GAA (2023) & 3\,nm & 7{,}78 & true QD channel \\
2\,nm (2025) & 2\,nm & 7{,}73 & $\approx$10 atomic layers \\
\bottomrule
\end{tabular}
\end{center}

\section{Biological Quantum Effects}

Living cells use quantum effects at 37\,°C. Three classes with
different conditions: Confinement quantization requires both
conditions; tunneling and spin coherence have their own criteria.
Condition~1 (geometric precision) applies to all three without exception.

\begin{center}
\small
\begin{tabular}{p{1.8cm}p{2.2cm}p{2.2cm}p{3cm}}
\toprule
\textbf{Effect} & \textbf{Cond.1: Geometry} & \textbf{Cond.2: $\Delta E > k_BT$} & \textbf{Mechanism} \\
\midrule
Photosynthesis (FMO) & \checkmark\ protein shell & borderline ($J/k_BT\approx0{,}7$) & exciton coherence, $J\approx19$\,meV \\
Bird navigation & \checkmark\ cryptochrome & no confinement & spin entanglement, relaxation time $>$ reaction time \\
Enzyme tunneling & \checkmark\ active site & no confinement & quantum tunneling through barrier \\
\bottomrule
\end{tabular}
\end{center}

The biological operating point at 37\,°C is not there because of a
thermal energy limit; $R_{\mathrm{bulk}}$ changes by less than 10\,\% between
0\,°C and 57\,°C. What collapses at 42--45\,°C
is the geometric precision: Proteins denature, the
torsion configuration collapses. 37\,°C is the point at which
biological resonators (protein cavities $\sim2$--4\,nm) are maximally chemically
active without losing their configuration.

\section{Influencing Spin: Five Methods}

Spin-up and spin-down are the absolute minimum of a quantum resonator:
two allowed orientations of a torsion configuration. Five
methods to rotate this orientation:

\textbf{Magnetic field (Zeeman effect).} The static field $B$ splits the
spin states by $\Delta E_Z = g^*\mu_BB$; a resonant
high-frequency pulse at the Larmor frequency rotates the spin in a targeted way. In the
FFGFT, the field establishes a preferred direction in the torsion field and
energetically breaks the symmetry between up and down.\status{E}

\textbf{Electric field (EDSR).} Electric fields couple via
spin-orbit coupling: The electron sees an
effective magnetic field in its rest frame, which rotates the spin. Advantage: Gate electrodes
generate fields localized to a few nanometers and switch in less than a
nanosecond. For integrated qubit arrays, this is the preferred
method.\status{E}

\textbf{Circularly polarized light.} $\sigma^+$ photons carry
angular momentum $+\hbar$ and transfer it upon resonant excitation to the
exciton and thus to the electron spin. Purely optical spin preparation
without an external field; readout via the polarization of the fluorescence light.
In the FFGFT: The photon carries its angular momentum as a topological
property of the electromagnetic field torsion; upon absorption,
it rotates the electron torsion.\status{E}

\textbf{Exchange interaction.} When two electrons in neighboring
quantum dots overlap spatially, their spins couple with $H_J =
J(t)\mathbf{S}_1\cdot\mathbf{S}_2$. By pulsing the gate field
(which controls the overlap and thus $J(t)$), one performs targeted
two-qubit operations --- the basis for spin qubit gates in
semiconductor systems (Loss-DiVincenzo, 1998).\status{E}

\textbf{Hyperfine interaction.} The electron interacts with the
nuclear spins of the $\sim10^4$--$10^6$ surrounding atoms. The fluctuating
nuclear spin bath normally destroys the spin polarization in $T_2^*\sim
10$--100\,ns --- a problem. Used in a controlled way, e.g., by
dynamical decoupling (rapid pulse sequences that average out nuclear spin fluctuations), $T_2$ is extended to microseconds. In phosphorus-in-silicon,
the nuclear spin itself is used as a highly stable qubit
($T_1 > 1$\,s at low temperatures).\status{E}

\section{Reading Out Spin: Five Ways}

\textbf{Spin-to-charge conversion.} The potential is chosen such
that only spin-down is energetically allowed to tunnel out of the quantum dot;
the tunneling electron generates a measurable current pulse in the
neighboring single-electron transistor. Accuracies $>99{,}5\,\%$ at
readout times below 1\,\textmu{}s; sufficient for error-corrected
quantum computers according to the surface code threshold.\status{E}

\textbf{Pauli blockade.} Two neighboring quantum dots, one electron each.
If one tries to load both into the same dot, this works for
antiparallel spins (both fit into the ground level), but not for parallel
spins (Pauli exclusion). Measurable as a conductance difference. In the FFGFT: Two
identical torsion configurations cannot occupy the same resonator,
not because of a prohibition, but because the same torsion winding in
the cavity does not allow a second stable configuration of the same kind. The
Pauli principle is a geometric statement.

\textbf{Faraday and Kerr rotation.} A linearly
polarized laser beam passes through the quantum dot without absorption and rotates its
polarization plane by $\theta_F\propto S_z$. This is a
quantum non-demolition measurement (QND); the spin state is unchanged.
Time resolution in the picosecond range; individual spins detectable.
Repeated measurements yield the same result.\status{E}

\textbf{Resonant fluorescence.} The photon is absorbed and excites the
transition --- invasive, the state is changed or destroyed afterwards.

\textbf{Imaging methods.} For ensembles: MOKE (magneto-optic
Kerr effect), time-resolved Faraday rotation, spin noise spectroscopy.

\section{Decoherence: The Limit of Superposition}

Two time scales describe the decay of quantum states.
$T_1$ (relaxation time) is the time until the system has fallen from the excited
state into the ground state. $T_2$ (coherence time) is the
time until the phase of the superposition is destroyed; always $T_2\le2T_1$.

\begin{center}
\small
\begin{tabular}{p{3.2cm}p{1.3cm}p{1.3cm}p{3.3cm}}
\toprule
\textbf{System} & \textbf{$T_1$} & \textbf{$T_2$} & \textbf{Main Mechanism} \\
\midrule
Electron spin GaAs & $\sim1$\,ms & $\sim10$\,ns & nuclear spin bath \\
Electron spin Si/SiGe & $>1$\,s & $\sim10$\,ms & ${}^{28}$Si (nuclear spin free) \\
Electron spin P:Si & $>1$\,s & $>100$\,ms & isotopic purification \\
Exciton CdSe QD & $\sim1$\,ns & $\sim10$\,ps & phonon scattering \\
Superconductor transmon & $\sim100$\,\textmu{}s & $\sim100$\,\textmu{}s & dielectric losses \\
\bottomrule
\end{tabular}
\end{center}

The comparison between exciton ($T_2\sim10$\,ps) and spin qubit in Si
($T_2\sim10$\,ms) yields nine orders of magnitude. With the Higgs coupling
parameter $\xi_{\mathrm{Higgs}}\approx1{,}038\times10^{-5}$, two hierarchy levels
bridge exactly $9{,}97$ orders of magnitude --- in agreement
with the measured ratio to $<5\,\%$.\status{S} The interpretation: The
exciton resonator couples directly to its environment (phonons of the same
level), the spin qubit to a distant bath (nuclear spins at level
$N\approx7$), thus two levels weaker.

Silicon enriched with ${}^{28}$Si eliminates the nuclear spin bath almost
completely. What remains is the intrinsic geometry of the
resonator. The long coherence time in ${}^{28}$Si is evidence that
geometric perfection (Condition~1) is the decisive one.\status{K}

\section{Superposition: Fragile and Never the Stable Ground State}

Here lies the actual conceptual core. The standard theory says:
A single electron is in a true superposition
$\alpha|\uparrow\rangle+\beta|\downarrow\rangle$ until measurement. On
every refrigerator there is a permanent magnet. Between these facts
lies a tension.

A permanent magnet is the macroscopic sum of spin orientations.
The exchange interaction forces neighboring spins into parallel
alignment; that is energetically more favorable. This collective alignment
forms Weiss domains. In the permanent magnet, the domains are permanently frozen by an external field. What the permanent magnet directly proves:
Spin has preferred, stable orientations; this stability
is macroscopically robust at room temperature over years; collective
superposition never occurs.

The standard theory explains this with decoherence: The interaction with
the environment destroys the superposition so quickly that it is practically
not observable. This is technically correct, but a subsequent
explanation: It begins with superposition as the ground state and then explains
why this ground state is always immediately destroyed.

The FFGFT begins the other way around: Stability is the ground state.
Superposition is the limited, fragile exception with measurable
time scales and a geometric cause.

\begin{center}
\small
\begin{tabular}{p{3.2cm}p{1.3cm}p{1.3cm}p{1.3cm}p{2.3cm}}
\toprule
\textbf{System} & \textbf{$T_1$} & \textbf{$T_2$} & \textbf{Temp.} & \textbf{Remark} \\
\midrule
Exciton CdSe QD & $\sim1$\,ns & $\sim10$\,ps & 4\,K & Superposition: picoseconds \\
Electron spin GaAs & $\sim1$\,ms & $\sim10$\,ns & mK & Superposition: nanoseconds \\
Electron spin Si & $>1$\,s & $\sim10$\,ms & mK & Superposition: milliseconds \\
Nuclear spin P:Si & $>1$\,s & $>100$\,ms & mK & best known system \\
\textbf{Permanent magnet} & \textbf{years} & --- & \textbf{300\,K} & \textbf{stability without cooling} \\
\bottomrule
\end{tabular}
\end{center}

The table shows: Not a single experimental system shows
superposition as a permanently stable state. Even in the best qubit
(nuclear spin P:Si at mK), the superposition decays after seconds. The
permanent magnet holds its stable spin state over years. Stability
is energetically preferred; superposition is metastable.

\textbf{The single electron.} The FFGFT position: A single electron
has a definite torsion configuration at every moment --- either
spin-up or spin-down. What is described as superposition $\alpha|\uparrow\rangle+
\beta|\downarrow\rangle$ is epistemic ignorance about this definite state,
not ontological simultaneity. The
proof comes from the QND measurement (Faraday rotation, Chapter~\ref{ch:qubit}):
A spin state can be measured repeatedly without changing it,
and always yields the same result. If the electron were truly in
superposition, a QND measurement would yield random
results on each repetition.\status{S}

\textbf{From single spin to permanent magnet.} The transition is, in the
FFGFT, not a puzzle, but the same mechanism at different scales:
Stability is geometrically grounded, superposition is a metastable
intermediate configuration, decoherence is the process that pushes the intermediate
torsion into one of the stable extremal configurations. In the
permanent magnet, $10^{23}$ torsion configurations have collectively
aligned to the same orientation and stabilize each other.
The macroscopic alignment is stable not despite, but because of the
geometric coupling.

\quelle{Doc.~173 (Quantum dot, March 2026), Doc.~174 (Spin and qubit states), Doc.~178 (Spin stability)}